\documentclass[a4paper]{article}

\usepackage{amsmath,amsfonts}

\setlength{\oddsidemargin}{0in}
\setlength{\textwidth}{6.5in}
\setlength{\topmargin}{0in}
\setlength{\headheight}{0in}
\setlength{\headsep}{0in}
\setlength{\textheight}{8.7in}
\usepackage[english]{babel}
\usepackage[utf8x]{inputenc}
\usepackage{amsmath}
\usepackage{graphicx}
\usepackage[colorinlistoftodos]{todonotes}

\title{Assignment 4}
\author{Frank Yang}

\begin{document}
\maketitle
\section{Homework list}

\begin{itemize}


\item Section  3.1:  18
\item Section 3.2:  6, 8, 13
\item Section 4.1:  4e, 8, 11, 19, 22ab, 24
\end{itemize}

\section{Solution}

\begin{enumerate}

\item[18]
Assume $a$ is an odd integer.By definition, there exists an integer $k$ such that $a=2k+1$. 
$$a^2-1=(2k+1)^2-1=(4k^2+4k+1)-1=4k^2+4k=4(k^2+k)$$
Let $p=k^2+k$. So $a^2-1=4p$, which means that $4$ is a factor of $a^2-1$.


\item[6]
Based on our previous proof, we grant that $\sqrt{2}$ is not rational.
We would prove all following statements  by contradiction.


(a) Suppose $-\sqrt{2}$ is rational. By definition, there exists two integers $m,n$ such that $-\sqrt{2}=\frac{m}{n}$. So $\sqrt{2}=\frac{-m}{n}$. Let $p=-m$. $\sqrt{2}=\frac{p}{n}$, where $p$ and $n$ are two integers. It contradicts our assumption that $\sqrt{2}$ is not rational. Thus, $-\sqrt{2} \notin \mathbb{Q}$.

(b) Suppose $1+\sqrt{2}$ is rational. By definition, there exists two integers $m,n$ such that $1+\sqrt{2}=\frac{m}{n}$. So $\sqrt{2}=\frac{m-n}{n}$. Let $p=m-n$. $\sqrt{2}=\frac{p}{n}$, where $p$ and $n$ are two integers. It contradicts our assumption that $\sqrt{2}$ is not rational. Thus, $1+\sqrt{2} \notin \mathbb{Q}$.

(c) Suppose $3+\sqrt{2}$ is rational. By definition, there exists two integers $m,n$ such that $3+\sqrt{2}=\frac{m}{n}$. So $\sqrt{2}=\frac{m-3n}{n}$. Let $p=m-3n$. $\sqrt{2}=\frac{p}{n}$, where $p$ and $n$ are two integers. It contradicts our assumption that $\sqrt{2}$ is not rational. Thus, $3+\sqrt{2} \notin \mathbb{Q}$.

(d) Suppose $r+\sqrt{2}$ is rational, where r is also a rational number. By definition, there exists two integers $m,n$ such that $r+\sqrt{2}=\frac{m}{n}$. Similarly, there exists two integers $t,k$ such that $r=\frac{t}{k}$ since $r$ itself is a rational number. Then,
$$\sqrt{2}=\frac{m}{n}-\frac{t}{k}$$
$$\sqrt{2}=\frac{km-nt}{nk}$$

Let $p=km-nt$ and $q=nk$. $\sqrt{2}=\frac{p}{q}$, where $p$ and $q$ are two integers. It contradicts our assumption that $\sqrt{2}$ is not rational. Thus, $r+\sqrt{2} \notin \mathbb{Q}$.



\item[8]
We give a proof by contradiction. 
Assume $\log_{2}(5)$ is rational. By definition,there exists two integers $m,n$ such that $\log_{2}(5)=\frac{m}{n}$. Then,
 $$2^{\log_{2}(5)}=2^{\frac{m}{n}}$$
 $$5=2^{\frac{m}{n}}$$
$$5^{n}=2^{m}$$
Apparently, $5^{n}$ is odd and $2^{m}$ is even. $5^{n}$ can never equal to $2^{m}$ for any two integers $m,n$. 
Thus,  $\log_{2}(5)$ is not rational.





\item[13]
We give a proof by contradiction. 
Assume $rx$ is rational, where r is a rational number but $\neq 0$ and x is not rational. By definition, there exists two integers $m,n$ such that $rx=\frac{m}{n}$. Similarly, there exists two integers $t,k$ such that $r=\frac{t}{k}$, since $r$ itself is a rational number. And since $r \neq 0$, $t$ is not $0$. 
Then,
 $$rx=\frac{tx}{k}=\frac{m}{n}$$
 $$txn=mk$$
 $$x=\frac{mk}{tn}$$
Let $p=mk$ and $q=tn$. $x=\frac{p}{q}$, where $p$ and $q$ are two integers. It contradicts our assumption that x is not a rational number. Thus, if $r$ is rational and $\neq 0$ and $x$ is irrational, then $rx$ is irrational. 





\item[4e]
$(A-C)=(A-B)\cup(B-C)$ is not true. 


Assume $A=\left\{{1,2}\right\}$ $B=\left\{{2,3}\right\}$ $C=\left\{{1,3}\right\}$


$(A-C)=\left\{{2}\right\}$ $(B-C)=\left\{{2}\right\}$ $(A-B)=\left\{{1}\right\}$ 

$(A-B)\cup(B-C)=\left\{{1,2}\right\} \neq (A-C)$


Salveage: $(A-C) \subseteq (A-B)\cup(B-C)$


Proof: Assume $x \in (A-C)$. Then, $x \in A$ and $x \notin C$. Either $x \in B$ or $x \notin B$. So, either $x \in B$ and $\notin C$, or $x \notin B$ and $x \in A$. So $x \in (B-C)$ or $x \in (A-B)$. So $x \in (A-B)\cup(B-C)$. Thus, $(A-C) \subseteq (A-B)\cup(B-C)$









\item[8]
(a) Assume $A \subseteq B$. So for all set $X$ that $X \subseteq A$, $X \subseteq B$. Since every element in $P(A)$ must be an subset of $A$, which means it must be an subset of $B$. And every subset of $B$ $\in P(B)$. So every element in $P(A)$ is also an element in $P(B)$. Thus, $P(A) \subseteq P(B)$.  


(b) By definition, for all element $X \in P(A \cap B)$, $X \subseteq (A \cap B)$. So, $X \subseteq A$ and $X \subseteq B$. Since $X \subseteq A$, $X \in P(A)$. Similarly, since $X \subseteq B$, $X \in P(B)$. So $X \in (P(A) \cap P(B)) $. Thus $P(A \cap B) \subseteq (P(A) \cap P(B))$. Since all previous steps are reversible. $(P(A) \cap P(B)) \subseteq P(A \cap B)$. Thus $P(A \cap B) = P(A) \cap P(B)$.


(c) $P(A \cup B)= P(A) \cup P(B)$


By definition, for all element $X \in P(A \cup B)$, $X \subseteq (A \cup B)$. So, $X \subseteq A$ or $X \subseteq B$. So $X \in P(A)$ or $X \in P(B)$. So $X \in (P(A) \cup P(B)) $. Thus $P(A \cup B) \subseteq (P(A) \cup P(B))$. Since all previous steps are reversible. $(P(A) \cup P(B)) \subseteq P(A \cup B)$. Thus $P(A \cup B) = P(A) \cup P(B)$.




\item[11]
On the last page. 
   
   
\item[19]
Proof by contrapositive:
Assume $B \neq C$. There exists an element $x$ that ($x$ in $B$ but not in $C$) or ($x$ in $C$ but not in $B$). Because $B$ and $C$ are symmetrical in this statement (swatching $B$ and $C$ will give us same proposition). Just assume $x$ in $B$ but not in $C$. For an element $a \in A$, the ordered pair $(a,x)$ will be in $A \times B$ but not in $A \times C$, since $x \notin C$. So $A \times B \neq A\times C$. Thus, if $A \times B= A\times C$, $B = C$.


\item[22a]
True.


Assume ordered pair $(p,q) \in A \times (B-C)$. So $p \in A$ and $q \in B$ and $q \notin C$. So $(p,q) \in A \times B$ and $(p,q) \notin A \times C$. So $(p,q) \in (A \times B - A\times C)$. Thus, $A \times (B-C) \subseteq A \times B - A\times C$. Since previous steps are all reversible, $A \times B - A\times C\subseteq A \times (B-C)$. So $A \times (B-C) = A \times B - A\times C$.   

\item[22b]
False.


Counterexample: 


Assume $A=\left\{{1}\right\}$


$B=\left\{{2}\right\}$


$U=\left\{{1,2}\right\}$


$A \times B =\left\{{(1,2)}\right\}$


$U^2=\left\{{(1,1),(1,2),(2,1),(2,2)}\right\}$


$(A \times B)^{c}=\left\{{(1,1),(2,1),(2,2)}\right\}$


$A^{c} \times B^{c}=\left\{{(2,1)}\right\}$


Thus, $(A \times B)^{c} \neq A^{c} \times B^{c}$.


\item[24]
(a)$\left\{{[0,1],[1,2],[2,3],[3,4],[4,5],[5,6],[6,7],[7,8],[8,9],[9,10],[10,11]}\right\}$


(b) $B_{1}=(0,2)$, $B_{10}=(0.9,1.1)$, $B_{100}=(0.99,1.01)$

\end{enumerate}
\end{document}